\chapter*{参考文献}\addcontentsline{toc}{chapter}{参考文献}
% Source PDF physical page 434; see audit/source_coverage.json.
1.建议同学们认真学习清华大学法学院张明楷教授的“两书一课”，真正的经典作品。

“\CJKunderwave{两书}”：（1）《张明楷刑法学讲义》，新星出版社2021年版（“得到”出品）。这是初阶教材，适合入门学习。（2）《刑法学》（第七版），法律出版社2026年版。这是著名的“太皇太后”，是进阶教材，适合专业学习。

“\CJKunderwave{一课}”，是指“得到”APP上的课程《刑法学100讲》，兼具趣味性、严谨性、权威性。

2.张明楷：《诈骗犯罪论》，法律出版社2021年版。

3.张明楷：《侵犯人身罪与侵犯财产罪》，北京大学出版社2021年版。

4.张明楷：《刑法的私垫》系列，北京大学出版社。

5.王钢：《刑法犯罪论核心知识与阶层式案例分析》，北京大学出版社2025年版。

6.王钢：《德国判例刑法》（分则），北京大学出版社2016年版。

7.李立众：《刑法一本通》（第十八版），法律出版社2025年版。

8.陈兴良主编：《案例刑法研究》，中国人民大学出版社2020年版。

9.周光权主编：《如何解答刑法题》，北京大学出版社2021年版。

10.周光权：《刑法总论》、《刑法各论》（第四版），中国人民大学出版社2021年版。

11.黎宏：《刑法学总论》、《刑法学各论》（第三版），法律出版社2025年版。

12.［日]山口厚：《刑法总论》（第3版），付立庆译，中国人民大学出版社2018年版。

13.[日]西田典之：《日本刑法总论》，刘明祥、王昭武译，中国人民大学出版社2007年版。

14.[日]大谷实：《刑法总论讲义》（新版第2版），黎宏译，中国人民大学出版社2008年版。

15.[德] Kindhauser，Strafrecht Allgemeiner Teil,6.Aufl.，Nomos，2013.

16.[德] Wessels/Beulke/Satzger,Strafrecht Allgemeiner Teil,43.Aufl.，C. F.Miller,2013.

17.[德] Rudolf Rengier,Strafrecht Allgemeiner Teil,4.Aufl.，Verlag C.H.Beck，2012.

18.[德] Claus Roxin,Strafrecht Allgemeiner Teil,Bd.I，Verlag C.H.Beck,2006.

19.最高人民法院主编的《刑事审判参考》。最高人民法院、最高人民检察院公布的指导性案例。

20.张明楷教授、王钢教授、李立众教授近年来的学术论文。
